\section{第~\ref{sec:recognition}~章　识别}

识别速度、最初几级的结构、响应的线性读出以及从时间连续性中学习，已在人和猴身上测量；把整条链归结为两种操作以及从视频中学习，已在本书的模型上检验，而对错觉的解释则从有充分依据到存在争议不等。

{\footnotesize
\setlength{\tabcolsep}{3pt}\renewcommand{\arraystretch}{1.15}
\begin{longtable}{>{\raggedright}p{0.5cm}>{\raggedright}p{4.0cm}>{\raggedright}p{5.6cm}>{\raggedright}p{2.3cm}>{\raggedright\arraybackslash}p{1.7cm}}
\toprule
序号 & 命题 & 实验及测量内容 & 文献 & 状态 \\
\midrule\endfirsthead
\toprule
序号 & 命题 & 实验及测量内容 & 文献 & 状态 \\
\midrule\endhead
1 & 平移时，逐像素与样板比较的正确率不比硬币好；一对级~\eqref{eq:scstage}能在任何位置认出图形 & \texttt{models/\allowbreak recognition\_models.py}，$24$点“视网膜”，$4000$次试验：翻转点比例为$0$、$0.1$、$0.2$时，样板法为$0.49$、$0.51$、$0.52$；$S$与$C$级对为$1.00$、$0.86$、$0.70$ & 本章推导 & 模型 \\
2 & 图片上的动物在$\sim150$\,ms内被认出，沿约十级走一遍，每级$10$--$20$\,ms & 人，对呈现$20$\,ms的照片判断“有无动物”：两种回答的诱发电位在$\sim150$\,ms后分开。猴：首次响应的时间逐级增加（V1最早，然后是V2和V4）。级数和“每级零个或一个脉冲”是由这些数字推出的 & \cite{Thorpe1996,Schmolesky1998} & 已测量 \\
3 & $S$级是对部件的“与”，$C$级是对位置取最大值（命题~\ref{prop:conveyor}） & 猫，初级视皮层：简单细胞响应特定位置、特定朝向的边缘，复杂细胞在更大区域内响应同一朝向。复杂细胞的细胞内记录：许多细胞对两根条纹的响应接近两个响应中较大的那个，平均而言最接近max运算。通过相互抑制取最大值见命题~\ref{prop:wta}，除以总活动见综述 & \cite{HubelWiesel1962,Lampl2004,Carandini2012} & 已测量 \\
4 & 沿链条向上，组合越来越大、越来越复杂，响应越来越不依赖位置 & 猴，把刺激简化为“关键特征”：V4和下颞叶皮层后部的细胞以较小的感受野响应形状、颜色和纹理的复杂组合；下颞叶皮层前部的感受野更大。模型中$S$与$C$的交替再现了这一图景 & \cite{KobatakeTanaka1994,Riesenhuber1999} & 已测量 \\
5 & 卷积网络重复了这种交替，并且对高层响应的预测最好；物体可以从神经元群体的频率中线性读出 & 为识别而训练、未对大脑做拟合的网络：顶层预测猴下颞叶皮层神经元的响应，中间层预测V4的响应。基于$\sim100$个随机选取细胞的活动、时间窗从$12.5$\,ms起的线性分类器，能在不同位置和大小下认出物体及其类别 & \cite{Fukushima1980,Yamins2014,Hung2005} & 已测量 \\
6 & 若资格迹不短于$\sim20$\,ms，带资格迹的赫布规则~\eqref{eq:traceinvariance}能从无标注视频中学会不变性 & 慢特征和资格迹规则的思想是理论。\texttt{models/\allowbreak recognition\_models.py}，两个$C$神经元，$16$个输入，运行$10$次：物体间停顿$0.3$\,s。$\tau_{\text{tr}}=5$\,ms时与图形的吻合平均为$0.52$，$10$\,ms时为$0.56$；$20$、$50$和$200$\,ms时所有运行均为$1.00$（$30$\,ms时十次中有一次出错） & \cite{Foldiak1991,WiskottSejnowski2002} & 模型 \\
7 & 大脑中的不变性是从时间的连续性中学到的 & 猴：在眼睛跳向视野某一位置时悄悄替换物体；下颞叶皮层神经元的位置不变性随替换而改变，一小时后即已显著。这是视觉学习的主要规则，则是假说 & \cite{LiDiCarlo2008} & 已测量 \\
8 & 运动：方向检测器，然后在大区域上使用同样的“与/或”对 & 兔视网膜中的方向检测器；猴MT区记录到的全部$110$个神经元都有方向选择性，对运动的响应强于V1 & \cite{Barlow1965,Albright1984} & 已测量 \\
9 & 高层向下发送预测，向上传送误差 & 该模型解释了初级视皮层感受野的非经典效应；个别观察与之相符（综述），但作为流水线的总体结构未经证明 & \cite{RaoBallard1999,Keller2018} & 假说 \\
10 & 困难图片需要反馈和绕几圈 & 猴：对“困难”图片，下颞叶皮层的判决晚$\sim30$\,ms形成；这些较晚的响应用前馈网络预测得差，用带反馈的或更深的网络预测得好 & \cite{Kar2019} & 已测量 \\
11 & 难以察觉的像素篡改能欺骗网络；人的视觉稳健得多，但并非无懈可击 & 在网络上：专门挑选的微小扰动改变答案；“熊猫$\to$长臂猿”的例子出自第二篇文献。对人，在短暂呈现时，能在网络间良好迁移的篡改会改变选择 & \cite{Szegedy2014,Goodfellow2015,Elsayed2018} & 已测量 \\
12 & 预期错觉：不可靠的输入让位于预期——暗淡的物体移动“更慢”，连衣裙被看成不同颜色（第~\ref{sec:illusions}~节） & 对比度较低的光栅看起来移动得较慢。对$1401$人的调查：连衣裙被稳定地看成三种样子，光照提示可切换颜色；在$\sim13\,000$名观察者中，起决定作用的是对光照的假定。预期慢速运动的贝叶斯模型解释了一系列运动错觉 & \cite{Thompson1982,LaferSousa2015,Wallisch2017,Weiss2002,Purves2011} & 假说 \\
13 & 增益调节：瀑布、倾斜和对比度；赫尔曼栅格不是邻域抑制 & 兔视网膜的方向检测器在长时间运动中降低频率，运动停止后降到背景以下。倾斜错觉可由除以邻居活动的模型再现。不改变视网膜输出神经元响应的栅格改动消除了暗斑——邻域抑制的解释被否定 & \cite{BarlowHill1963,Schwartz2009,SchillerCarvey2005} & 已测量 \\
14 & 延迟补偿：运动物体看起来在闪光的前方 & 错觉已测量。心理物理学结果既与运动外推矛盾，也与延迟差矛盾；有人提出事后解释——依据闪光后$\sim80$\,ms内的事件。争论尚无定论 & \cite{Nijhawan1994,EaglemanSejnowski2000} & 假说 \\
15 & 静止的“蛇”在旋转：延迟差~\eqref{eq:latency}被方向检测器读出 & 错觉在无色时也存在；相邻元素对单独就能产生它；猴皮层中具有方向选择性的神经元对这些静止的元素对给出方向性响应。在人身上，旋转由眼睛的微跳视和眨眼触发 & \cite{Conway2005,OteroMillan2012} & 已测量 \\
16 & 双稳图形：相互抑制、疲劳和噪声；切换主要由噪声引起 & 竞争神经元群体的模型：其中的切换由\emph{噪声}引起，没有噪声就没有切换；它解释了切换频率随刺激强度增加。疲劳在这里作用较小 & \cite{MorenoBote2007} & 假说 \\
17 & 一部分错觉是机器所学任务的结果 & 训练来预测视频下一帧的网络“看到”静止圆环在旋转，而在对照图片上看不到；用于图像去噪和锐化的网络像人一样受颜色和亮度错觉影响 & \cite{Watanabe2018,GomezVilla2020} & 已测量 \\
\bottomrule
\end{longtable}}
